\section{Analysis}

The analyses in this section are separated from the primary performance comparisons because they are intended to characterize heterogeneity and generate hypotheses, rather than to establish additional confirmatory claims. Unless stated otherwise, subgroup results are descriptive and are not interpreted as biological mechanisms.

\subsection{Task Heterogeneity}

Performance varies substantially across the 157 human tissue--HLA tasks. On the standard benchmark, the worst-10 mean AUROC is 0.6834 despite an overall mean of 0.8448. The ten lowest task AUROCs are concentrated in lung, blood, and lymphoid labels, ranging from 0.6510 for lung--HLA-B*49:01 to 0.7056 for lymphoid--HLA-B*13:02. This pattern may reflect biological similarity, heterogeneous composition, annotation, sampling, or study effects; the present data cannot distinguish these explanations.

Table~\ref{tab:human-tissue-extremes} reports the five lowest and five highest tissue-macro summaries rather than selecting isolated tasks. Single-task tissues remain unstable subgroup estimates and are shown descriptively only.

\begin{table}[tbp]
\centering
\caption{Descriptive human tissue-level extremes on the 157-task standard benchmark. Means are taken across available HLA tasks within each tissue.}
\label{tab:human-tissue-extremes}
\footnotesize
\setlength{\tabcolsep}{2pt}
\begin{tabular}{lrrrlrrr}
\toprule
Lower tissue & Tasks & AUROC & AUPRC & Higher tissue & Tasks & AUROC & AUPRC \\ 
\midrule
Blood & 22 & 0.7618 & 0.7587 & Testis & 1 & 0.9573 & 0.9597 \\ 
Lymphoid & 34 & 0.7914 & 0.7744 & Tongue & 2 & 0.9574 & 0.9535 \\ 
Uterine cervix & 7 & 0.8021 & 0.7797 & Uterus & 3 & 0.9617 & 0.9598 \\ 
Brain & 6 & 0.8085 & 0.8012 & Adrenal gland & 2 & 0.9778 & 0.9768 \\ 
Lung & 19 & 0.8135 & 0.8081 & Umbilical cord blood & 1 & 0.9791 & 0.9805 \\ 
\bottomrule
\end{tabular}
\end{table}

Task difficulty is only weakly associated with the number of fitting rows. Spearman correlation between training size and standard AUROC is \(\rho=-0.109\) (\(p=0.176\)); for peptide-disjoint AUROC it is \(\rho=0.135\) (\(p=0.093\)). Neither association is significant at 0.05, so more rows do not guarantee a higher task AUROC in this benchmark.

For a separate descriptive view by HLA locus, Table~\ref{tab:human-hla-locus} groups the original fixed-test and strict OOF AUROCs. HLA-A has the highest fixed-test mean and HLA-C the lowest; HLA-C also shows the smallest descriptive decrease.

\begin{table}[tbp]
\centering
\caption{Descriptive human AUROC by HLA locus. Difference is strict OOF minus the original fixed test; because the held-out pools differ, it is not an estimate of the peptide-overlap effect.}
\label{tab:human-hla-locus}
\small
\begin{tabular}{lrrrr}
\toprule
HLA locus & Tasks & Original fixed & Strict OOF & Difference \\ 
\midrule
HLA-A & 55 & 0.8688 & 0.7824 & -0.0864 \\ 
HLA-B & 79 & 0.8380 & 0.7507 & -0.0873 \\ 
HLA-C & 23 & 0.8106 & 0.7738 & -0.0369 \\ 
\bottomrule
\end{tabular}
\end{table}

These locus summaries contain different alleles, tissues, training sizes, and peptide-component structures and therefore are not controlled locus effects. Differences are calculated from unrounded task means, so subtracting the displayed values can differ by 0.0001. The matched standard-OOF comparison appropriate for estimating the peptide-separation effect appears in RQ4 and Figure~\ref{fig:appendix-gap}.

Standard fixed-test and strict-OOF AUROCs correlate moderately across tasks (Spearman \(\rho=0.569\), \(p<10^{-14}\)). Their unmatched median difference is -0.0652, with 16 increases and 141 decreases. Six of the ten largest descriptive decreases involve HLA-A*11:01 and two involve HLA-B*49:01, motivating allele- and component-aware follow-up without establishing intrinsic allele effects.

\subsection{Branch Complementarity}

We analyze the global auxiliary and HLA-specific branches using held-out predictions from the three-seed auxiliary dual-branch experiment. Within each seed and task, the Pearson correlation between branch probabilities has a mean of 0.706 and a median of 0.718, with a range from 0.312 to 0.945. The branches are therefore related but not interchangeable: some tasks receive nearly redundant orderings, whereas others show substantial disagreement.

After averaging each branch's task AUROC across seeds, the global branch outperforms the HLA-specific branch on 81 tasks, while the HLA-specific branch wins on 76 tasks; there are no exact ties. This near-even division is notable because the global branch has a higher average AUROC overall. The HLA-specific branch can still provide useful task-level information even when its mean standalone performance is lower.

The rank-fused model outperforms the better of its two component branches in 260 of 471 seed--task comparisons. Its median improvement over the stronger branch is 0.0017, while the mean improvement is 0.0003. Fusion is therefore beneficial for a majority of comparisons but not universally so. Moreover, a simple disagreement measure, \(1-\operatorname{corr}(s_g,s_h)\), is not significantly associated with fusion gain (Spearman \(\rho=-0.062\), \(p=0.181\)). Large disagreement alone is insufficient to predict complementarity: disagreement may represent either useful independent signal or branch-specific error.

These results support fixed fusion as a robust aggregate rule, but they do not imply that every task benefits from both branches.

\subsection{Error Analysis}

Available task-level results point to several recurring sources of difficulty. First, low-performing tasks are concentrated in tissue categories such as blood and lymphoid, where closely related cell populations and broad sampling may produce overlapping peptide evidence. Second, sparse tasks can have unstable estimates even though training size alone does not explain performance. Third, peptides reported across several tissues can become difficult pseudo-negatives when an unobserved target-tissue record reflects incomplete measurement rather than true absence. Finally, study and batch structure may induce both easy shortcuts and systematic errors when experimental coverage is uneven.
